\chapter{Matchings}

\section{Basic definitions and augmenting paths}

\begin{notedefinition}
Let $G=(V,E)$ be a graph.
\begin{itemize}
  \item A \emph{matching} is a set $M\subseteq E$ whose edges have pairwise
        disjoint endpoint sets.
  \item A matching is \emph{maximal} if no edge can be added to it, and
        \emph{maximum} if it has largest possible cardinality.
  \item A vertex incident with an edge of $M$ is \emph{saturated} by $M$;
        otherwise it is \emph{unsaturated} or \emph{exposed}.
  \item A \emph{perfect matching} saturates every vertex.
\end{itemize}
The maximum size of a matching is denoted $\nu(G)$.
\end{notedefinition}

\begin{notedefinition}
Let $M$ be a matching.  A path is \emph{$M$-alternating} if its edges
alternate between $M$ and $E\setminus M$.  It is \emph{$M$-augmenting} if
both endpoints are unsaturated by $M$.  Every augmenting path has odd length.
\end{notedefinition}

\begin{notedefinition}
For sets $A$ and $B$, their \emph{symmetric difference} is
\[
  A\triangle B=(A\setminus B)\cup(B\setminus A).
\]
\end{notedefinition}
\begin{notetheorem}
If $G$ is connected and $\Delta(G)\le2$, then $G$ is a path or a cycle.
\end{notetheorem}
\begin{proof}
If every vertex has degree $2$, start at any vertex and follow unused edges.
Because the graph is finite and no vertex permits branching, the first repeat
closes a cycle.  Connectedness then forces that cycle to contain every vertex.

Otherwise $G$ has a vertex $v$ of degree at most $1$.  Starting from $v$ and
extending a path as far as possible produces no repeated vertex, since a repeat
would create a cycle and leave no available edge through which the path could
have entered or left it.  If a vertex lay outside the resulting maximal path,
connectedness would give an edge from the path to the outside and hence a
further extension.  Thus the path contains all vertices.
\end{proof}

\begin{notetheorem}[Berge's Theorem]
A matching $M$ is maximum if and only if there is no $M$-augmenting path.
\end{notetheorem}
\begin{proof}
If $P$ is $M$-augmenting, then
\[
  M':=M\triangle E(P)
\]
is a matching and $|M'|=|M|+1$, so $M$ is not maximum.

Conversely, suppose a matching $N$ satisfies $|N|>|M|$.  In the spanning
subgraph with edge set $M\triangle N$, every vertex has degree at most $2$.
Its components are therefore paths and even cycles whose edges alternate
between $M$ and $N$.  Since $N$ has more edges overall, some path component
contains one more edge of $N$ than of $M$.  That path starts and ends with
$N$-edges, so both endpoints are unsaturated by $M$; it is an
$M$-augmenting path.
\end{proof}

\begin{noteremark}
Berge's theorem turns maximum matching into a search problem: repeatedly find
an augmenting path and take the symmetric difference.  In bipartite graphs,
this search can also be implemented as a network-flow computation.
\end{noteremark}
\section{Hall's theorem and its consequences}

\begin{notedefinition}
A graph is \emph{$k$-regular} if every vertex has degree $k$.  It is
\emph{regular} if it is $k$-regular for some $k\ge0$.
\end{notedefinition}

\begin{noteproposition}
If $k\ge1$, every $k$-regular bipartite graph $G=(X,Y;E)$ satisfies
$|X|=|Y|$.
\end{noteproposition}
\begin{proof}
Counting edges from both sides gives
\[
  k|X|=\sum_{x\in X}\deg(x)=|E|
  =\sum_{y\in Y}\deg(y)=k|Y|.
\]
\end{proof}

\begin{notetheorem}
Every $k$-regular bipartite graph with $k\ge1$ has a perfect matching.
\end{notetheorem}
\begin{proof}
This follows from Hall's theorem below.  For $S\subseteq X$, count the edges
between $S$ and $N(S)$:
\[
  k|S|
  =e(S,N(S))
  \le k|N(S)|.
\]
Thus $|S|\le|N(S)|$, so there is a matching saturating $X$.  Since
$|X|=|Y|$, it is perfect.
\end{proof}
\begin{notedefinition}
For $S\subseteq V(G)$, the \emph{neighbourhood} of $S$ is
\[
  N(S):=\{v\in V(G):uv\in E(G)\text{ for some }u\in S\}.
\]
A matching is \emph{$S$-saturating} if it saturates every vertex of $S$.
\end{notedefinition}

\begin{notetheorem}[Hall's Marriage Theorem]
Let $G=(X,Y;E)$ be bipartite.  There is a matching saturating $X$ if and only
if
\[
  |S|\le|N(S)|\qquad\text{for every }S\subseteq X.
\]
\end{notetheorem}
\begin{proof}
Necessity is immediate: distinct vertices of $S$ must be matched to distinct
vertices of $N(S)$.

For sufficiency, induct on $|X|$.  The result is trivial when $|X|=1$.
Assume $|X|>1$.

First suppose every nonempty proper subset $S\subsetneq X$ satisfies the
strict inequality $|N(S)|\ge|S|+1$.  Choose an edge $xy$.  In
$G-\{x,y\}$, for every $S\subseteq X\setminus\{x\}$,
\[
  |N_{G-\{x,y\}}(S)|\ge |N_G(S)|-1\ge |S|.
\]
Induction gives a matching saturating $X\setminus\{x\}$, and adding $xy$
finishes.

Otherwise choose a nonempty proper $S\subsetneq X$ with $|N(S)|=|S|$.
Hall's condition holds in the induced bipartite graph $G[S,N(S)]$, so by
induction it has a matching saturating $S$.  It also holds in
$G[X\setminus S,Y\setminus N(S)]$: for $R\subseteq X\setminus S$,
\[
  |N(R)\setminus N(S)|
  =|N(R\cup S)|-|N(S)|
  \ge |R\cup S|-|S|=|R|.
\]
Induction gives a matching saturating $X\setminus S$.  The two matchings are
vertex-disjoint, and their union saturates $X$.
\end{proof}
\begin{noteapplication}
Theorem~6.4 is an immediate application of Hall's theorem: every
$k$-regular bipartite graph with $k\ge1$ has a perfect matching.
\end{noteapplication}

\begin{notedefinition}
Let $S_1,\dots,S_n$ be finite sets.  A \emph{system of distinct
representatives} is a choice of elements $x_i\in S_i$ such that
$x_1,\dots,x_n$ are pairwise distinct.
\end{notedefinition}

\begin{notetheorem}[Hall's theorem for set systems]
The family $(S_1,\dots,S_n)$ has a system of distinct representatives if and
only if
\[
  |I|\le\left|\bigcup_{i\in I}S_i\right|
  \qquad\text{for every }I\subseteq\{1,\dots,n\}.
\]
\end{notetheorem}
\begin{proof}
Build a bipartite graph whose left vertices are the indices $1,\dots,n$ and
whose right vertices are the elements appearing in the sets, joining $i$ to
$x$ exactly when $x\in S_i$.  A matching saturating the left side is exactly
a system of distinct representatives, and Hall's condition becomes the
displayed inequality.
\end{proof}
\begin{notedefinition}
An $m\times n$ \emph{Latin rectangle} is an $m\times n$ matrix with entries
from an $n$-element symbol set such that no symbol is repeated in any row or
column.  A Latin square is an $n\times n$ Latin rectangle.
\end{notedefinition}

\begin{notetheorem}
Every $m\times n$ Latin rectangle can be extended to an $n\times n$ Latin
square by adding $n-m$ rows.
\end{notetheorem}
\begin{proof}
To add one row, form a bipartite graph whose left side consists of the $n$
columns and whose right side consists of the $n$ symbols.  Join column $j$ to
symbol $a$ exactly when $a$ has not yet appeared in column $j$.  Every column
is missing $n-m$ symbols.  Every symbol has appeared once in each of the $m$
rows, hence in $m$ distinct columns, so it too has degree $n-m$.  The graph is
regular bipartite and therefore has a perfect matching.  Put in each column
the symbol assigned by this matching.  This creates a new valid row.  Repeat
until $n$ rows have been obtained.
\end{proof}
\section{Doubly stochastic matrices}

\begin{notedefinition}
A square matrix is \emph{doubly stochastic} if all entries are nonnegative and
every row sum and every column sum is $1$.  A \emph{permutation matrix} is a
$0$--$1$ matrix with exactly one $1$ in every row and every column.
\end{notedefinition}

\begin{notedefinition}
A \emph{convex combination} of objects $v_1,\dots,v_k$ in a real vector space
is an expression
\[
  \lambda_1v_1+\cdots+\lambda_kv_k,
  \qquad
  \lambda_i\ge0,
  \quad
  \sum_{i=1}^k\lambda_i=1.
\]
\end{notedefinition}

\begin{notetheorem}
Every convex combination of permutation matrices is doubly stochastic.
\end{notetheorem}
\begin{proof}
Nonnegativity is preserved, and each row or column sum is the corresponding
convex combination of $1$'s, hence equals $1$.
\end{proof}

\begin{notetheorem}[Birkhoff--von Neumann Theorem]
Every doubly stochastic matrix is a convex combination of permutation
matrices.
\end{notetheorem}
\begin{proof}
Let $A=(a_{ij})$ be doubly stochastic.  Form the support graph with a row
vertex $r_i$, a column vertex $c_j$, and an edge $r_ic_j$ whenever
$a_{ij}>0$.  For any set $S$ of row vertices, all mass in those rows lies in
$N(S)$, so
\[
  |S|=\sum_{r_i\in S}\sum_j a_{ij}
  \le\sum_{c_j\in N(S)}\sum_i a_{ij}=|N(S)|.
\]
Hall's theorem gives a perfect matching, hence a permutation matrix $P$ whose
$1$-entries occur at positive entries of $A$.

Let $\mu$ be the minimum selected entry.  If $\mu=1$, then the row and column
sum conditions force $A=P$.  Otherwise
\[
  B:=\frac{A-\mu P}{1-\mu}
\]
is again doubly stochastic and has fewer positive entries than $A$.  Induction
on the number of positive entries writes $B$ as a convex combination of
permutation matrices, and
\[
  A=\mu P+(1-\mu)B
\]
gives the required representation.
\end{proof}
\section{Vertex covers and stable matchings}

\begin{notedefinition}
A \emph{vertex cover} is a set $C\subseteq V(G)$ incident with every edge.
The minimum size of a vertex cover is denoted $\tau(G)$.
\end{notedefinition}

\begin{notetheorem}[K\H{o}nig--Egerv\'ary Theorem]
For every bipartite graph $G$,
\[
  \nu(G)=\tau(G).
\]
\end{notetheorem}
\begin{proof}
Every vertex cover meets every edge of a matching, so $\nu(G)\le\tau(G)$.
For the reverse inequality, let $G=(X,Y;E)$ and let $M$ be a maximum matching.
Let $U\subseteq X$ be the vertices unsaturated by $M$, and let $Z$ be the set
of vertices reachable from $U$ by $M$-alternating paths that begin with an
edge outside $M$.  Put $Z_X=Z\cap X$ and $Z_Y=Z\cap Y$.

There is no augmenting path, so every vertex of $Z_Y$ is matched.  Its matched
partner lies in $Z_X\setminus U$, and this gives a bijection
$Z_Y\leftrightarrow Z_X\setminus U$.  Hence
\[
  |Z_Y|=|Z_X|-|U|.
\]
Now
\[
  C:=(X\setminus Z_X)\cup Z_Y
\]
is a vertex cover.  Indeed, an edge from a vertex of $Z_X$ to a vertex outside
$Z_Y$ would extend an alternating reachability path.  Finally,
\[
  |C|=|X|-|Z_X|+|Z_Y|=|X|-|U|=|M|.
\]
Thus $\tau(G)\le|C|=\nu(G)$.
\end{proof}

\begin{noteremark}
Under the standard matching-to-flow reduction, the minimum vertex cover in
K\H{o}nig's theorem corresponds to a minimum cut.
\end{noteremark}
\begin{notedefinition}
Let $G=K_{n,n}$ with bipartition $(X,Y)$.  Suppose every vertex has a strict
ranking of the opposite part.  A perfect matching $M$ is \emph{stable} if
there is no unmatched pair $xy$ such that $x$ prefers $y$ to its partner in
$M$ and $y$ prefers $x$ to its partner in $M$.
\end{notedefinition}

\begin{notetheorem}
Every such preference instance has a stable perfect matching.
\end{notetheorem}

\begin{notealgorithm}[Gale--Shapley]
Initially everyone is unmatched.  While some $x\in X$ is unmatched and has
not proposed to every vertex of $Y$, let $x$ propose to the most-preferred
vertex $y$ to whom it has not yet proposed.  The vertex $y$ keeps the better
of its current proposer and $x$, rejecting the other.  When no proposal is
possible, return the held pairs.
\end{notealgorithm}
\begin{proof}[Correctness]
No ordered pair is considered twice, so the algorithm stops after at most
$n^2$ proposals.  Once a vertex of $Y$ receives a proposal, it remains matched
and can only replace its partner by a more-preferred one.  If some $x\in X$
were unmatched at termination, it would have proposed to every $y\in Y$; all
$n$ vertices of $Y$ would then be holding proposals from the remaining
$n-1$ vertices of $X$, which is impossible.  Thus the returned matching is
perfect.

Suppose an unmatched pair $xy$ blocked the returned matching.  Since $x$
prefers $y$ to its final partner, $x$ proposed to $y$ earlier.  At that time
$y$ rejected $x$, immediately or later, for a partner it preferred.  Because
$y$ only trades upward, it prefers its final partner to $x$, a contradiction.
\end{proof}
\begin{noteremark}
The bipartite hypothesis is essential for this formulation.  Stable matching
in general graphs is a different problem and requires a different algorithm.
\end{noteremark}
\section{The Tutte--Berge formula}

\begin{notedefinition}
For $S\subseteq V(G)$, let $\odd(G-S)$ be the number of odd-order components
of $G-S$, and define
\[
  \deficiency(S):=\odd(G-S)-|S|.
\]
\end{notedefinition}

\begin{notelemma}[deficiency lemma]
Every finite graph $G$ has a set $S\subseteq V(G)$ and a maximum matching $M$
such that
\begin{itemize}
  \item every even component of $G-S$ has a perfect matching;
  \item every odd component of $G-S$ has a matching missing exactly one
        vertex;
  \item $M$ matches the vertices of $S$ to distinct odd components of
        $G-S$.
\end{itemize}
Consequently, $M$ misses exactly $\odd(G-S)-|S|$ vertices.
\end{notelemma}
\begin{proof}[Proof sketch]
Choose $S$ so that $\deficiency(S)$ is maximum and, subject to this, $|S|$ is
maximum.  A standard alternating-path argument shows that every even
component of $G-S$ is factorable, every odd component is factor-critical, and
the bipartite incidence graph between $S$ and the odd components satisfies
Hall's condition.  Match $S$ into distinct odd components, use a near-perfect
matching in each selected odd component that leaves the attachment vertex
free, and use perfect or near-perfect matchings in all remaining components.
If any of these assertions failed, an alternating-path exchange would either
increase the deficiency or enlarge $S$ without decreasing it, contradicting
the choice of $S$.
\end{proof}

\begin{notetheorem}[Tutte--Berge Formula]
If $G$ has $n\ge1$ vertices, then
\[
  \nu(G)
  =\frac12\min_{S\subseteq V(G)}
    \bigl(n+|S|-\odd(G-S)\bigr)
  =\frac12\left(n-
    \max_{S\subseteq V(G)}\deficiency(S)\right).
\]
\end{notetheorem}
\begin{proof}
Let $M$ be any matching and $S\subseteq V(G)$.  At most $|S|$ odd components
of $G-S$ can send an $M$-edge to $S$.  Every remaining odd component contains
an exposed vertex.  Hence
\[
  n-2|M|\ge \odd(G-S)-|S|,
\]
or equivalently
\[
  |M|\le\frac12\bigl(n+|S|-\odd(G-S)\bigr).
\]
This proves the upper bound for every $S$.  The deficiency lemma supplies a
set $S$ and a matching that misses exactly $\odd(G-S)-|S|$ vertices, so
its size is
\[
  \frac12\bigl(n+|S|-\odd(G-S)\bigr).
\]
The upper bound is therefore attained.
\end{proof}

\begin{notecorollary}[Tutte's 1-Factor Theorem]
A graph $G$ has a perfect matching if and only if
\[
  \odd(G-S)\le|S|
  \qquad\text{for every }S\subseteq V(G).
\]
\end{notecorollary}
\begin{proof}
By the Tutte--Berge formula, $\nu(G)=|V(G)|/2$ exactly when
$\deficiency(S)\le0$ for every $S\subseteq V(G)$, which is precisely the
displayed condition.
\end{proof}
\section{Exercises and further consequences}

\begin{noteexercise}
Let $G$ be bipartite.  Then $G$ has a perfect matching if and only if
\[
  \alpha(G)=\frac{|V(G)|}{2},
\]
where $\alpha(G)$ is the maximum size of an independent set.
\end{noteexercise}
\begin{proof}
If $M$ is perfect, an independent set contains at most one endpoint of each
edge of $M$, so $\alpha(G)\le|V|/2$.  Either bipartition class is independent,
so equality holds.

Conversely, the complement of a vertex cover is independent and conversely,
so $\alpha(G)+\tau(G)=|V|$.  If $\alpha(G)=|V|/2$, then
$\tau(G)=|V|/2$.  K\H{o}nig's theorem gives
$\nu(G)=\tau(G)=|V|/2$, hence a perfect matching.
\end{proof}

\begin{noteexercise}
Every bipartite graph $G$ with at least one edge has a matching of size at
least
\[
  \frac{|E(G)|}{\Delta(G)}.
\]
\end{noteexercise}
\begin{proof}
Let $C$ be a minimum vertex cover.  Since every vertex of $C$ is incident with
at most $\Delta(G)$ edges,
\[
  |E(G)|\le\Delta(G)|C|.
\]
By K\H{o}nig's theorem, $|C|=\nu(G)$.
\end{proof}
\begin{noteexercise}
Let $T$ be a tree.
\begin{enumerate}
  \item If $T$ has a perfect matching, then it is unique.
  \item The tree $T$ has a perfect matching if and only if
        \[
          \odd(T-v)=1\qquad\text{for every }v\in V(T).
        \]
\end{enumerate}
\end{noteexercise}
\begin{proof}
For (1), if $M$ and $M'$ were distinct perfect matchings, every component of
$M\triangle M'$ would be an even cycle, impossible in a tree.

For (2), suppose $M$ is perfect and remove $v$.  Exactly one component of
$T-v$ contains the vertex matched to $v$; that component has odd order, while
every other component is perfectly matched internally and has even order.
Thus $\odd(T-v)=1$.

Conversely, assume the displayed condition and induct on $|V(T)|$.  The case
$|V(T)|=2$ is immediate.  Choose a leaf $x$ and let $y$ be its neighbour.
The singleton $\{x\}$ is an odd component of $T-y$; because
$\odd(T-y)=1$, every other component $C_1,\dots,C_k$ of $T-y$ has even
order.

For $u\in V(C_i)$, all vertices outside $C_i$ lie in one connected part of
$T-u$ attached through $y$, and that outside part has even order.  Joining it
to the relevant component of $C_i-u$ therefore does not change the parity of
that component.  Consequently
\[
  \odd(C_i-u)=\odd(T-u)=1
  \qquad (u\in V(C_i)).
\]
Each $C_i$ is smaller than $T$, so induction gives a perfect matching in every
$C_i$.  Together with the edge $xy$, these matchings form a perfect matching
of $T$.
\end{proof}
\begin{noteexercise}[Petersen's Theorem]
Every bridgeless $3$-regular graph has a perfect matching.
\end{noteexercise}
\begin{proof}
By Tutte's theorem it suffices to prove
$\odd(G-S)\le|S|$ for every $S\subseteq V(G)$.  Let $C$ be an odd component
of $G-S$.  Degree counting gives
\[
  3|V(C)|=2|E(C)|+|[V(C),S]|,
\]
so the cut $[V(C),S]$ has odd size.  It cannot have size $1$, since that edge
would be a bridge.  Hence every odd component sends at least three edges to
$S$.  Therefore
\[
  3\odd(G-S)
  \le |[S,V(G)\setminus S]|
  \le 3|S|,
\]
and Tutte's condition follows.
\end{proof}
